\documentclass[10pt,a4paper]{amsart}
\usepackage[margin=19mm]{geometry}
\usepackage{amsmath,amssymb,mathtools}
\usepackage[scr=rsfs,cal=euler]{mathalfa}
\usepackage{xfrac}
\usepackage[unicode,hidelinks,bookmarks=false]{hyperref}
\newcommand{\ie}{i.e.\ }
\newcommand{\cf}{cf.\ }
\newcommand{\resp}{respectively\ }
\newcommand{\Subsection}{\subsection}
\providecommand{\nobreakdash}{\nobreak\hbox{-}}
\setlength{\emergencystretch}{2em}
\multlinegap0pt
\usepackage{bm}
\usepackage{bbm}
\usepackage{dsfont}
\usepackage{xspace}
\usepackage{cancel}
\usepackage{mathtools}
\usepackage{amsthm}
\makeatletter
\@ifundefined{theorem}{\newtheorem{theorem}{Theorem}}{}
\@ifundefined{lemma}{\newtheorem{lemma}[theorem]{Lemma}}{}
\@ifundefined{proposition}{\newtheorem{proposition}[theorem]{Proposition}}{}
\makeatother
\DeclareMathOperator{\AFa}{AF_\beta}
\DeclareMathOperator{\AF}{AF}
\DeclareMathOperator{\ALFm}{ALF^{--}}
\DeclareMathOperator{\ALFp}{ALF^+}
\DeclareMathOperator{\Ae}{AE}
\newcommand{\mR} {\mathfrak R}
\title[Gravitational instantons and harmonic maps]{Gravitational instantons and harmonic maps}
\author{Mingyang Li, Song Sun}
\date{Source: arXiv:2507.15284v1 (CC BY 4.0)}
\begin{document}
\maketitle

\noindent\emph{Source and licence.} Gravitational instantons and harmonic maps. Version arXiv:2507.15284v1 (CC BY 4.0), distributed under the Creative Commons Attribution 4.0 International licence, as recorded in the arXiv licence metadata for this exact version and shown on the abstract page https://arxiv.org/abs/2507.15284v1. Full rights notice: licenses/arxiv-2507.15284-CC-BY-4.0.txt.\par\smallskip

\noindent\textit{Editorial scope.} Counted unit: the Proposition whose source label is \texttt{prop:appendix 1} in the article (source0.tex lines 3987--3991, source label \texttt{prop:appendix 1}), delivered together with the article's own complete original proof of it (source0.tex lines 3992--4082, one \texttt{proof} environment of 91 lines). No mathematics is written, repaired or re-derived: statement and proof are the article's own text line for line. Required context retained uncounted: thm:C11 regularity (lines 1333--1341); lem:Lipschitz metric (lines 1389--1391). Every definition, display and label of those lines is retained unchanged. Omitted from this package and not redistributed: 1--1332, 1342--1388, 1392--3986, 4083--4255. Nothing inside a retained excerpt is omitted or altered: every retained body is delivered exactly as the article's lines are written, so no omission receipt of any kind accompanies this package. Citations: the retained excerpts carry no citation command at all, so no bibliography entry of the article is transcribed and no reference is redistributed. Reference adaptation: none was needed and none was made; every reference inside the delivered unit resolves inside it, so no unresolved reference or citation is produced. Printed slips kept literal: no printed word, symbol, hypothesis or number of the counted statement or of its proof was altered. Layout and class: the article's own class is \texttt{amsart} with packages including inputenc, fontenc, lmodern, babel, microtype, amsmath, amssymb, amsfonts, amsthm, cases, mathtools, accents, mathrsfs, aliascnt; the wrapper is the lane's pinned \texttt{amsart} editorial wrapper at 10pt on A4, which restates the theorem-like environments the retained text needs and adds the article's own macro definitions that the retained text needs (\texttt{\textbackslash AF}, \texttt{\textbackslash AFa}, \texttt{\textbackslash ALFm}, \texttt{\textbackslash ALFp}, \texttt{\textbackslash Ae}, \texttt{\textbackslash mR}), with extra wrapper packages (bm, bbm, dsfont, xspace, cancel, mathtools). The delivered numbering is the unit's own. Assets: no retained span contains a figure environment, a graphics-inclusion command, a PDF-page inclusion, a table or a TikZ picture; no image file is copied into this package, the package declares no asset dependency and no image file is redistributed with it. Licence: version arXiv:2507.15284v1 of this article is distributed under the Creative Commons Attribution 4.0 International licence, as recorded in the arXiv licence metadata for this exact version and shown on the abstract page https://arxiv.org/abs/2507.15284v1. A full rights notice is delivered at licenses/arxiv-2507.15284-CC-BY-4.0.txt. Characters: every retained excerpt is pure ASCII, so the delivered engine, XeLaTeX with its default Latin Modern text font, reports no missing character.
\begin{proposition}[$C^{1,\alpha}$ regularity]\label{thm:C11 regularity}
	We have the following uniform bounds on $B_{\frac{1}{2}}(0)$:
	\begin{align*}
		&|\nabla^{k}v|\leq CC_0\rho^{2-k},\text{ for $k=0,1,2,3$},\\
		&\|u\|_{C^{1,\alpha}}\leq CC_0,\text{ for any $0<\alpha<1$},\\
        &|\nabla^ku|\leq CC_0\rho^{1-k}\text{ for any $k=1,2,3$}.
	\end{align*}
    Here the constant $C$ depends on $\alpha$ but is independent of $C_0$.
\end{proposition}

\begin{lemma}\label{lem:Lipschitz metric}
	The Ricci-flat metric $g$ extends to a Lipschitz (i.e. $C^{0, 1}$) metric on $V$.
\end{lemma}

\begin{proposition}\label{prop:appendix 1}

    
    Given a rod structure $\mR$ of  Asymptotic Type  $\sharp\in\{\Ae,\ALFm,\ALFp, \AFa\}$ and a harmonic map $\Phi$ strongly tamed by $\mR$, with the normalized augmentation $\nu$ and the lattice $\Lambda_\mR$, the associated smooth Ricci-flat end $(E,g_{\Lambda_\mR})$ also has  Asymptotic Type  $\sharp$ with decay rate $\delta=1$.
\end{proposition}

\begin{proof}
    We prove the proposition for the $\AF_0$ case and the other cases are similar. Over the end $E$, which is the completion of $(\mathbb{H}^\circ\setminus B_N(0))\times\mathbb{R}^2/\Lambda_\mR$, there is the model metric $g^{\AF_0}$. It suffices to find a diffeomorphism $\varpi:E\to E$, such that $\varpi_*g_{\Lambda_\mR}$ is smooth, and 
    $$|\nabla_{g^{\AF_0}}^k(\varpi_*g_{\Lambda_{\mR}}-g^{\AF_0})|_{g^{\AF_0}}=O(\mathfrak{r}^{-k-1}).$$

    We construct the diffeomorphism on each annuli $A_{2^{j},2^{j+2}}$ with $j\in\mathbb{N}$ large, then patch them together. By the modified rescaling $\Phi_{2^j}=D_{\sqrt{2^j}}.\varphi_{2^{-j}}^*\Phi$, we scale the harmonic map $\Phi$ on $A_{2^{j-1},2^{j+3}}$ to $\Phi_{2^j}$ on $A_{1/2,8}$ which satisfies $d(\Phi_{2^j},\Phi^{\AF_0})\leq C2^{-j}$. Set
    $$\Phi_{2^j}=\frac{1}{\rho}\left(\begin{matrix}
        \rho^2e^{u_j}+e^{-u_j}v_j^2 & e^{-u_j}v_j \\ e^{-u_j}v_j & e^{-u_j} 
    \end{matrix}\right),\quad \Phi=\frac{1}{\rho}\left(\begin{matrix}
        \rho^2e^{u}+e^{-u}v^2 & e^{-u}v \\ e^{-u}v & e^{-u} 
    \end{matrix}\right)$$
    over $A_{1/2,8}$ and $A_{2^{j-1},2^{j+3}}$ as before, then $u_j=\varphi_{2^{-j}}^*u$ and $v_j=2^{-j}\varphi_{2^{-j}}^*v$. The uniform local regularity Proposition \ref{thm:C11 regularity} gives that on $A_{1/2,8}$
    \begin{equation}\label{eq:C1a scale to A12 appendix}
        \|u_j\|_{C^{1,\alpha}}\leq C2^{-j},\quad |\nabla^kv_j|\leq C_k2^{-j}\rho^{2-k}
    \end{equation}
    for $0<\alpha<1$ and $k=0,1,2$. 


    \vspace{10pt}
    \noindent\textbf{Step (1). Construction of local gauge in each annulus.}
    \vspace{10pt}


    

    Now, under the augmentation $\varphi_{2^{-j}}^*\nu$ on $A_{1/2,8}$, with the lattice $\Lambda_{\mR}$, we get  the Ricci-flat metric $h$ associated to $\Phi_{2^j}$ defined on the completion of $A_{1/2,8}\times\mathbb{R}^2/\Lambda_{\mR}$ (denoted by $E_{1/2,8}$).  Define the sector region
    $$Sec_\tau:=\{\tau<\theta<\pi-\tau\}\subset\mathbb{H}^\circ.$$
    For a fixed small $\tau$, over $A_{1/2,8}\cap Sec_\tau$, the harmonic map equation is non-degenerate. Hence, by \eqref{eq:C1a scale to A12 appendix} on $A_{1/2,8}\cap Sec_\tau$ standard elliptic estimates give the improvement 
    \begin{equation}
        \|u_j\|_{C^{k,\alpha}}\leq C_k2^{-j},\quad \|v_j\|_{C^{k,\alpha}}\leq C_k2^{-j}
    \end{equation}
    for any $k\in\mathbb{N}$. Therefore, on the region $\pi^{-1}(A_{1/2,8}\cap Sec_\tau)\subset E_{1/2,8}$, under the coordinate $(x_1,x_2,x_3,x_4):=(r\cos\phi_1,r\sin\phi_1,z,\phi_2)$, we already have
    $$|\nabla_{g^{\AF_0}}^k(h-g^{\AF_0})|_{g^{\AF_0}}\leq \frac{C_k}{2^j}.$$
    
    
    
    As for the region in $A_{1,4}$ near the rod, we cover them by balls $B_t\subset A_{1/2,8}$ with fixed radius $1/4$. Then we cover the near rod region in $E_{1,4}:=\pi^{-1}(A_{1,4})$ by simply-connected open sets $U_{t,s}\subset E_{1/2,8}$ with smooth controlled boundary over $B_t$.  Introduce the coordinate
    $$x_1:=r\cos\phi_1,\quad x_2:=r\sin\phi_1,\quad x_3:=z,\quad x_4:=\phi_2$$
    on $U_{t,s}$. On each $U_{t,s}$, we solve for the harmonic coordinates
    \begin{equation}
        {h}^{kl}\frac{\partial^2}{\partial x_k\partial x_l}\overline{x}_j+{h}^{kl}{\Gamma}^r_{kl}\frac{\partial}{\partial x_r}\overline{x}_j=0,\quad \text{with }\overline{x}_j=x_j\text{ along }\partial U_{t,s}.\label{eq:harmonic coordinates for h appendix}
    \end{equation}
    Since we know the metric $h$ is $C^{0,1}$ under $(x_1,x_2,x_3,x_4)$ by Lemma \ref{lem:Lipschitz metric}, the harmonic coordinates are $W^{2,p}$ for any large $p$, hence is $C^{1,\alpha}$. Our uniform local regularity Proposition \ref{thm:C11 regularity} gives that under $(x_1,x_2,x_3,x_4)$
    \begin{equation}\label{eq:C11 bound for h appendix}
        |h_{kl}-g^{\AF_0}_{kl}|+|\Gamma^r_{kl}|\leq\frac{C}{2^j}.
    \end{equation}
    The elliptic estimate based on \eqref{eq:harmonic coordinates for h appendix} gives $\|\overline{x}_j-x_j\|_{W^{2,p}}\leq C2^{-j}$ for any large $p$, which implies $\|\overline{x}_j-x_j\|_{C^{1,\alpha}}\leq C2^{-j}$ under $(x_1,x_2,x_3,x_4)$. The Ricci-flat metric $h$ is smooth under $(\overline{x}_1,\overline{x}_2,\overline{x}_3,\overline{x}_4)$. Define the following map where we use the coordinates $(x_1,x_2,x_3,x_4)$ for $U_{t,s}$
    $$\zeta_{t,s}:U_{t,s}\to U_{t,s},$$
    $$x\mapsto \zeta_{t,s}(x):=(\overline{x}_1,\overline{x}_2,\overline{x}_3,\overline{x}_4).$$
    Shrinking the domain for $\zeta_{t,s}$ slightly, the map is a $C^{1,\alpha}$ diffeomorphism into its image when $j$ is large. Now the metric $(\zeta_{t,s})_*h$ under $(x_1,x_2,x_3,x_4)$ is smooth. The bound $\|\overline{x}_j-x_j\|_{C^{1,\alpha}}\leq C2^{-j}$ and \eqref{eq:C11 bound for h appendix} imply
    $$\|(\zeta_{t,s})_*h-g^{\AF_0}_{kl}\|_{C^{0,\alpha}}\leq \frac{C}{2^j}.$$
    Since the Ricci-flat equation is elliptic in harmonic coordinates, starting from the above, it is standard to show $\|(\zeta_{t,s})_*h-g^{\AF_0}_{kl}\|_{C^{k,\alpha}}\leq C_k2^{-j}$.




    \vspace{10pt}
    \noindent\textbf{Step (2). Equivariant patching to global gauge in each annulus.}
    \vspace{10pt}



    In the next we need to patch local diffeomorphisms on $\pi^{-1}(A_{1/2,8}\cap Sec_\tau)$ and $U_{t,s}$ together to construct a global equivariant diffeomorphism on $E_{1,4}$. Denote the identity diffeomorphism on $U_{0,0}:=\pi^{-1}(A_{1,4}\cap Sec_\tau)$ as $\zeta_{0,0}$. For the open cover of $E_{1,4}$ by $U_{t,s}$, we take standard cut-off functions $0\leq \chi_{t,s}\leq1$ that
    \begin{itemize}
        \item is smooth under the harmonic coordinate for $U_{t,s}$;
        \item has support in $U_{t,s}$; 
        \item equals one in a slightly smaller open subset of $U_{t,s}$ and $\sum\chi_{t,s}\equiv1$.
    \end{itemize}
    Consider the center-of-mass function
    $$\widetilde{\mathfrak{C}}_j:E_{1/2,8}\times E_{1/2,8}\to\mathbb{R}_+,$$
    $$(x,y)\mapsto \sum_{t,s}\chi_{t,s}(x)\cdot d_{h}(\zeta_{t,s}(x),y)^2.$$
    The function is almost invariant under the $\mathbb{T}$ action, in the sense that for $a\in \mathbb{T}$, $$\|\widetilde{\mathfrak{C}}_j(a\cdot x,a\cdot y)-\widetilde{\mathfrak{C}}_j(x,y)\|_{C^{k,\alpha}}\leq C_k2^{-j},$$ where the $C^{k,\alpha}$ norm is measured under the metric $h$. Denote by $\mathfrak{C}_j$ the average of  $\widetilde{\mathfrak{C}}_j$ over the $\mathbb T$ orbits.
    Now for any $x\in E_{1,4}$, we can minimize the function $y\mapsto \widetilde{\mathfrak{E}}_j(x,y)$. The minimizing point must take place at $y$  close to $x$, hence by convexity, the minimizing point is unique. Denote the minimizing point as $\mathfrak{z}_j(x)$. Then the map 
    $$\mathfrak{z}_j:E_{1,4}\to E_{1/2,8}$$
    is a $\mathbb{T}$-equivariant diffeomorphism into its image, and the push-forward metric satisfies
    $$\|(\mathfrak{z}_j)_*h-g^{\AF_0}_{kl}\|_{C^{k,\alpha}}\leq\frac{C_k}{2^j}.$$



    To recover the diffeomorphism on the annulus $A_{2^k,2^{j+2}}$, we can scale $A_{1,4}$ back to $A_{2^j,2^{j+2}}$, where the metric $h$ is rescaled to $2^{2j}h$. Because of the modified rescaling for $\Phi$, after further taking the quotient of $2^{2j}h$ by the $\mathbb{Z}_{2^j}$ action along $\phi_2$, we exactly recover the metric $g_{\Lambda_{\mR}}$ on $\pi^{-1}(A_{2^j,2^{j+2}})$. Since $\mathfrak{z}_j$ is $\mathbb{T}$-equivariant, it induces a $\mathbb{T}$-equivariant diffeomorphism $\overline{\mathfrak{z}}_j$ on $\pi^{-1}(A_{2^j,2^{j+2}})$ smoothing $g_{\Lambda_{\mR}}$ with the desired regularity.




    \vspace{10pt}
    \noindent\textbf{Step (3). Equivariant patching of adjacent annuli.}
    \vspace{10pt}



    Finally, we need to patch the diffeomorphisms $\overline{\mathfrak{z}}_j$ and $\overline{\mathfrak{z}}_{j+1}$ on the overlap of adjacent annuli $A_{2^j,2^{j+2}}$ and $A_{2^{j+1},2^{j+3}}$. Recall for the construction of $\mathfrak{z}_{j+1}$, we need to consider the modified rescaling $\Phi_{2^{j+1}}$ on $A_{1/2,8}$ and repeat the construction in \textbf{Step (1)}, which is equivalent to consider the modified rescaling $\Phi_{2^j}$ on $A_{1,16}$. Therefore, it suffices to patch the $\mathbb{T}$-equivariant diffeomorphisms $\mathfrak{z}_j$ on $A_{1,4}$ and $\mathfrak{z}_{j+1}$ on $A_{2,8}$. One just needs to perform the center-of-mass construction again. Hence we have obtained the desired global diffeomorphism $\varpi$.
\end{proof}

\end{document}
